
\documentclass{article}
\usepackage{colortbl}
\usepackage{makecell}
\usepackage{multirow}
\usepackage{supertabular}

\begin{document}

\newcounter{utterance}

\centering \large Interaction Transcript for game `ifeval', experiment `ifeval', episode 190 with Qwen3.5{-}4B.
\vspace{24pt}

{ \footnotesize  \setcounter{utterance}{1}
\setlength{\tabcolsep}{0pt}
\begin{supertabular}{c@{$\;$}|p{.15\linewidth}@{}p{.15\linewidth}p{.15\linewidth}p{.15\linewidth}p{.15\linewidth}p{.15\linewidth}}
    \# & \multicolumn{2}{c}{Player} && \multicolumn{2}{c}{Game Master} \\
    \hline

    \theutterance \stepcounter{utterance}  
    & & & \multicolumn{4}{p{0.6\linewidth}}{
        \cellcolor[rgb]{0.9,0.9,0.9}{
            \makecell[{{p{\linewidth}}}]{
                \texttt{\tiny{[P1$\langle$GM]}}
                \texttt{Are the weather conditions in the Arctic very cold most of the year? Do not say 'yes' or 'no' throughout your entire response.} \\
            }
        }
    }
     \\ \\

    \theutterance \stepcounter{utterance}  
    & \multicolumn{4}{p{0.6\linewidth}}{
        \cellcolor[rgb]{0.561,0.737,0.561}{
            \makecell[{{p{\linewidth}}}]{
                \texttt{\tiny{[P1$\rangle$GM]}}
                \texttt{The Arctic is indeed characterized by extremely low temperatures for a significant portion of the year, particularly during the winter months when the sun remains below the horizon for weeks, leading to prolonged darkness and freezing conditions. However, describing the climate as cold "most of the year" requires some nuance, as the region experiences a distinct seasonal cycle.} \\
\\ 
\texttt{During the summer, the Arctic undergoes a phenomenon known as the midnight sun, where the sun stays above the horizon for 24 hours a day. While temperatures can still drop below freezing in the interior regions and during the early and late parts of the summer, coastal areas and the ocean surface often remain relatively mild, sometimes even warming above freezing. This creates a contrast where the land may freeze while the sea remains liquid.} \\
\\ 
\texttt{Furthermore, the definition of "cold" can vary depending on the specific location within the Arctic. The polar regions, such as the North Pole itself, experience more consistent freezing temperatures year{-}round, whereas the Arctic Ocean and surrounding islands like Svalbard or northern Norway see much milder summers compared to the deep winter lows. Therefore, while the Arctic is undeniably a frozen landscape for a large part of the annual cycle, the summer months offer a period of relative warmth that prevents the entire year from being uniformly and intensely cold.} \\
            }
        }
    }
    & & \\ \\

    \theutterance \stepcounter{utterance}  
    & & & \multicolumn{2}{p{0.3\linewidth}}{
        \cellcolor[rgb]{0.9,0.9,0.9}{
            \makecell[{{p{\linewidth}}}]{
                \texttt{\tiny{[GM$|$GM]}}
                {'keywords:forbidden_words': {'forbidden_words': ['yes', 'no']}}
            }
        }
    }
    & & \\ \\

    \theutterance \stepcounter{utterance}  
    & & & \multicolumn{2}{p{0.3\linewidth}}{
        \cellcolor[rgb]{0.9,0.9,0.9}{
            \makecell[{{p{\linewidth}}}]{
                \texttt{\tiny{[GM$|$GM]}}
                \texttt{game\_result = WIN} \\
            }
        }
    }
    & & \\ \\

\end{supertabular}
}

\end{document}
